% Source: https://artofproblemsolving.com/community/c3h1452065_useful_log_identity
% Stable identity: aops-topic-1452065-q01
\textbf{(a)} By the power rule $\log_a(x^k)=k\log_a x$ with $k=\log_a y$,
\[
\log_a\!\left(x^{\log_a y}\right)=\log_a y\cdot\log_a x .
\]

\textbf{(b)} By (a) with the roles of $x$ and $y$ exchanged, $\log_a\!\left(y^{\log_a x}\right)=\log_a x\cdot\log_a y$ as well. So $x^{\log_a y}$ and $y^{\log_a x}$ have the same logarithm to base $a$, and since $\log_a$ is one-to-one,
\[
x^{\log_a y}=y^{\log_a x}.
\]
For example, $9^{\log_3 5}=5^{\log_3 9}=5^2=\boxed{25}$, and $8^{\log_2 7}=7^{\log_2 8}$, so $8^{\log_2 7}\cdot7^{-\log_2 8}=\boxed{1}$.

There is no product rule that turns $\log_a x\cdot\log_a y$ into a single logarithm of a simple expression in $x$ and $y$ alone; the identity in (a) (or the change-of-base formula $\log_a x\cdot\log_x y=\log_a y$) is the useful form.
